\documentclass[11pt,a4paper]{article}
% XeLaTeX / Tectonic. TeX Live font files; no OS font lookup or shell escape.
\usepackage{fontspec}
\setmainfont{cmunrm.otf}[BoldFont=cmunbx.otf,ItalicFont=cmunti.otf,BoldItalicFont=cmunbi.otf]
\usepackage[a4paper,margin=20mm]{geometry}
\usepackage{amsmath,amssymb,booktabs,tabularx,graphicx,xcolor}
\usepackage[font=small,labelfont=bf]{caption}
\usepackage[numbers,sort&compress]{natbib}
\usepackage{microtype}
\usepackage[section]{placeins}
\usepackage[colorlinks=true,linkcolor=blue!45!black,citecolor=blue!45!black,urlcolor=blue!45!black]{hyperref}
\setlength{\parindent}{1em}
\setlength{\parskip}{2pt}
\setlength{\textfloatsep}{13pt plus 2pt minus 2pt}
\setlength{\floatsep}{10pt plus 2pt minus 2pt}
\setlength{\emergencystretch}{2em}
\renewcommand{\topfraction}{0.90}
\renewcommand{\bottomfraction}{0.85}
\renewcommand{\textfraction}{0.08}
\renewcommand{\floatpagefraction}{0.75}
\newcommand{\norm}[1]{\frac{#1}{\lVert #1\rVert_2}}
\newcommand{\softplus}{\operatorname{softplus}}
\newcommand{\LSE}{\operatorname{LSE}}
\newcommand{\pending}[1]{\textcolor{black!65}{\textit{#1}}}

% Public source revision; availability must be verified separately.
\newcommand{\artifactroot}{https://github.com/iCodeCraft/reaction-lens-systems/blob/fdd417c93432495cd031d14772175d492a690ca5}
\newcommand{\repositorynotice}{Code and data: \url{https://github.com/iCodeCraft/reaction-lens-systems}.}
\newcommand{\evidence}[2]{\href{\artifactroot/#1}{#2}}

\newcommand{\SampleCount}{\evidence{NUMBERS.csv}{64}}
\newcommand{\RepeatCount}{\evidence{NUMBERS.csv}{3}}
\newcommand{\MemoryBatch}{\evidence{NUMBERS.csv}{16}}
\newcommand{\MemorySpans}{\evidence{NUMBERS.csv}{14}}
\newcommand{\ContractChecks}{\evidence{NUMBERS.csv}{768}}
\newcommand{\ScoreComparisons}{\evidence{NUMBERS.csv}{10545024}}
\newcommand{\MaxError}{\evidence{NUMBERS.csv}{0.0}}
\newcommand{\DecisionFlips}{\evidence{NUMBERS.csv}{0}}
\newcommand{\SmallBlock}{\evidence{NUMBERS.csv}{128}}
\newcommand{\SmallMedian}{\evidence{NUMBERS.csv}{1.676}}
\newcommand{\SmallPninetyfive}{\evidence{NUMBERS.csv}{1.827}}
\newcommand{\SmallMemory}{\evidence{NUMBERS.csv}{412.53}}
\newcommand{\DefaultBlock}{\evidence{NUMBERS.csv}{1024}}
\newcommand{\DefaultMedian}{\evidence{NUMBERS.csv}{0.642}}
\newcommand{\DefaultPninetyfive}{\evidence{NUMBERS.csv}{0.701}}
\newcommand{\DefaultMemory}{\evidence{NUMBERS.csv}{414.17}}
\newcommand{\FullBlock}{\evidence{NUMBERS.csv}{15692}}
\newcommand{\FullMedian}{\evidence{NUMBERS.csv}{0.656}}
\newcommand{\FullPninetyfive}{\evidence{NUMBERS.csv}{0.764}}
\newcommand{\FullMemory}{\evidence{NUMBERS.csv}{444.05}}
\newcommand{\EndRequests}{\evidence{NUMBERS.csv}{15}}
\newcommand{\EndExamples}{\evidence{NUMBERS.csv}{5}}
\newcommand{\EndMedian}{\evidence{NUMBERS.csv}{268.90}}
\newcommand{\EndPninetyfive}{\evidence{NUMBERS.csv}{393.54}}
\newcommand{\SetupSeconds}{\evidence{NUMBERS.csv}{4.87}}
\newcommand{\EndMemory}{\evidence{NUMBERS.csv}{1149.22}}
\newcommand{\EncodingShare}{\evidence{NUMBERS.csv}{96.33}}
\newcommand{\PublicationCount}{\evidence{NUMBERS.csv}{2377}}
\newcommand{\CatalogCount}{\evidence{NUMBERS.csv}{15692}}
\newcommand{\KnownLinks}{\evidence{NUMBERS.csv}{4094}}
\newcommand{\MeasuredMAP}{\evidence{NUMBERS.csv}{0.4770}}
\newcommand{\MeasuredHit}{\evidence{NUMBERS.csv}{0.3925}}
\newcommand{\MeasuredRecall}{\evidence{NUMBERS.csv}{0.9016}}
\newcommand{\MetricError}{\evidence{NUMBERS.csv}{0.0}}

\usepackage[english]{babel}
\title{Reaction Lens Systems: An Experience Report on Verifiable Full-Catalog Scientific Retrieval}
\author{Imran Gadzhiev\\Independent Researcher}
\date{5 October 2026}
\begin{document}
\maketitle
\begin{abstract}
Scientific retrieval software must preserve specified scoring behavior while its
candidate collection and execution strategy change. We report an engineering
evaluation of Reaction Lens Systems, an existing system that ranks Reactome reactions from
English scientific text. We derive candidate-transformation contracts from its
pointwise global/local scorer and check its blocked implementation against an
unblocked reference. On sampled development inputs, \ContractChecks{} checks
covering \ScoreComparisons{} score comparisons produce maximum observed logit
difference \MaxError{} and \DecisionFlips{} fixed-threshold decision changes.
For a catalog of \CatalogCount{} reactions, median cached-feature scoring time is
\DefaultMedian{} ms with the default block size, while median complete local
prediction time is \EndMedian{} ms on a separate convenience sample. An isolated
installation also reproduces the checkpoint's recorded development ranking metrics.
The contribution is an artifact-backed account of execution contracts, component
costs and reproducibility boundaries in one retrieval system; it is not a new
learning algorithm or a claim of general retrieval superiority.
\end{abstract}
{\small\noindent\repositorynotice}
\section{Problem and contribution}
A researcher inspecting a scientific abstract may need to identify relevant
biological reactions in a curated catalog. Reaction Lens Systems supplies ranked proposals
for human review. It does not generate reactions or establish that a proposed
association is biologically correct. Its engineering problem is maintaining an
inspectable relationship between input text, catalog entries and predictions.
Changing the order of candidates or the scoring block size should not silently
change the score assigned to an unchanged article--reaction pair.

This report studies an existing implementation rather than introducing a new
retrieval service. Its contributions are an explicit execution contract with
metamorphic checks, a controlled comparison of blocked and unblocked exact
scoring, and a package that separates regeneration of recorded tables from fresh
model inference. These results are intended to help maintainers distinguish
changes in ranking caused by new competitors from unintended changes in pairwise
scores, and to identify which component dominates local request cost.
We do not claim that candidate independence or blocked matrix operations are novel.

\section{System and related work}
Reactome supplies the biological reaction catalog \citep{reactome2026}.
SPECTER and SPECTER2 provide citation-informed scientific document representations
\citep{cohan2020specter,singh2023scirepeval}. Reaction Lens Systems freezes its encoder and
uses a learned global/local head. We evaluate a fixed checkpoint to study
execution equivalence, component costs and inference reproducibility.

The processing path is English title/abstract, deterministic sentence spans,
frozen text encoding, projected reaction vectors, exact blocked scoring, then
ranked and threshold-selected output. Reaction vectors are prepared once during
initialization. The Python API, CLI and local HTTP service share the prediction
implementation. The implementation rejects simultaneous calls on one instance;
we evaluate local sequential calls, not multi-user throughput. Bundle loading
verifies listed file hashes and shapes. These checks detect mismatches against a
manifest; they do not authenticate a maliciously replaced manifest or guarantee
semantic compatibility of arbitrary user-supplied encoders.

Metamorphic testing checks relations between transformed inputs and outputs when
individual expected outputs are difficult to specify \citep{segura2016}.
We instantiate this established approach for the candidate axis of the scorer.
Agreement with an unblocked implementation checks execution equivalence, not
biological correctness. Our reference shares the scoring implementation, so it
cannot expose all common-mode defects. No fault-detection superiority is claimed.

\section{Execution contracts}
Let $g$ and $s_j$ be the normalized projections of the global article and
valid sentence features through the shared learned article projection, and $r$
the normalized output of the separate learned reaction projection. The fixed head computes
\begin{equation}
 z(x,r)=a\left[g^\top r+w\max\left(0,\max_j s_j^\top r-n\right)\right]+b,
 \qquad \widehat{y}(x,r)=\mathbf{1}[z(x,r)\geq\tau].
\end{equation}
Here the learned scale $a$, local weight $w$, null $n$, bias $b$, and development-fitted
threshold $\tau$ are held fixed. The implementation uses
$a=\exp(\min(\ell,\log 100))$, $w=\operatorname{softplus}(u)$ and
$n=\tanh(v)$ for learned scalars $\ell,u,v$. Padding is masked before the
sentence maximum. No reduction is taken over competing reactions.
Consequently, in real arithmetic, permuting candidates permutes outputs; removing
candidates preserves retained scores; and appending candidates preserves previous
scores. Fixed-threshold decisions inherit this property. Rankings and top-$k$
membership do not: new candidates can outrank existing ones. Updating descriptions,
model weights, preprocessing or the threshold also falls outside this contract.

For batch size $B$, sentence count $S$, candidate count $R$ and projected width
$P$, local similarities require $O(BSRP)$ arithmetic. Candidate blocks of width
$K$ bound the principal local-similarity intermediate to $O(BSK)$ instead of
$O(BSR)$. The final score matrix still uses $O(BR)$ storage, and the prepared
reaction bank uses $O(RP)$ storage. Blocking is an execution tradeoff, not a
retrieval shortlist or a reduction in the number of scored candidates.

\section{Evaluation protocol}
The protocol is included as \texttt{PROTOCOL.md}. Experiments use the fixed
structured-description checkpoint and its checksummed cached development pack.
All \CatalogCount{} catalog reactions are scored. We sample \SampleCount{} publications
without replacement using NumPy seed 20261005 and run CPU FP32 with two PyTorch
threads. The environment is macOS 26.6.1 on arm64, Python 3.12.14, PyTorch 2.14.1
and NumPy 2.5.3. The sandbox did not expose the CPU brand; these results should
not be treated as a portable hardware performance specification.

Each block configuration receives five warmup queries and \RepeatCount{} timed repetitions
per sampled publication. Timings exclude encoding, bundle loading, initial
reaction projection and result serialization. Configurations run sequentially in
fixed order, so thermal and background-load effects remain possible. Sample p95
values are descriptive quantiles of repeated measurements, not confidence bounds.

On the first repetition we check candidate permutation, a half-catalog subset,
appending 128 duplicate reaction vectors, and an unblocked reference. The
predeclared absolute logit tolerance is $10^{-5}$. Duplicate additions exercise
candidate-axis execution; they do not evaluate new biological descriptions.
Threshold checks use the existing release threshold without recalibration.

Memory is measured separately in a fresh process for each configuration using
the same padded batch of \MemoryBatch{} sampled publications, with maximum sentence count \MemorySpans{}.
We record process high-water resident memory including imports and loaded feature
arrays. These are not incremental tensor allocations, and each configuration has
only one memory trial. Processes used for this comparison do not execute the
unblocked reference or the candidate-transformation probes.

Complete prediction measurements use \EndExamples{} existing released English
comparison examples and three repetitions after one warmup. Their selection was
for demonstration, so they are a convenience sample. Prediction includes text
encoding, scoring, ranking and payload construction. JSON serialization is recorded
separately; HTTP/network costs are excluded. Model setup is timed separately.

\section{Results}
\subsection{Candidate and execution equivalence}
The checks produce \ScoreComparisons{} scalar comparisons across
\ContractChecks{} query/configuration/scenario combinations. The maximum observed
absolute logit error is \MaxError{}, with \DecisionFlips{} fixed-threshold flips.
These measurements support the intended contract on this checkpoint, platform and
sample. They are not a universal bitwise-equivalence proof across backends.
Evidence: \evidence{results/table3_contracts.csv}{contract observations}, derived from all individual probes
in \texttt{results/system-benchmark.json}. Regenerate with \texttt{make contracts}.

\subsection{Blocking and local execution cost}
\begin{table}[ht]
\centering
\begin{tabular}{rrrr}
\toprule
Block width & Median (ms) & p95 (ms) & Peak RSS (MiB)\\
\midrule
\SmallBlock & \SmallMedian & \SmallPninetyfive & \SmallMemory\\
\DefaultBlock & \DefaultMedian & \DefaultPninetyfive & \DefaultMemory\\
\FullBlock & \FullMedian & \FullPninetyfive & \FullMemory\\
\bottomrule
\end{tabular}
\caption{\raggedright Cached-feature single-publication scoring latency and separate padded-batch
process memory. These are different workloads. Every timing row scores the full
catalog. Evidence: \evidence{results/table1_scoring.csv}{block comparison}; individual timings and memory
observations are in \texttt{results/system-benchmark.json}. Regenerate with \texttt{make table1}.}
\end{table}
The smallest block incurs greater observed latency. The default and unblocked
configurations have similar medians; this sample does not establish a significant
speed difference between them. The blocked configurations have lower observed
process peak RSS in the separate batch probe, consistent with the intermediate
storage bound. Memory differences should not be generalized from a single trial.

The \EndRequests{} complete local predictions have median \EndMedian{} ms and p95
\EndPninetyfive{} ms. Encoding accounts for \EncodingShare\% of aggregate measured
prediction time. Setup takes \SetupSeconds{} s, and process peak RSS reaches
\EndMemory{} MiB. Evidence: \evidence{results/table2_end_to_end.csv}{component timings} and its raw
source. Regenerate with \texttt{make timings}. For these examples, improving only head execution offers limited benefit
to total latency because encoding dominates. This observation does not measure
the effect of encoder caching, alternative encoders or approximate retrieval.

\subsection{Reproduction boundary}
An isolated dependency environment recomputes predictions from the downloaded
checkpoint and cached features for \PublicationCount{} development publications,
\KnownLinks{} known links and the full catalog. MAP is \MeasuredMAP{}, Hit@1 is
\MeasuredHit{}, and micro Recall@100 is \MeasuredRecall{}. Maximum absolute
difference from the stored release metrics is \MetricError{}.
Evidence: \evidence{results/fresh-evaluation.json}{fresh evaluation}. Recorded links are incomplete,
and development data were used during model development. Recall@100 is not
top-one accuracy or evidence of biological correctness.

For publication $i$, average precision sums the precision at each relevant
rank over the full catalog and divides by the number of distinct recorded
reaction labels for that publication, including labels absent from the catalog.
MAP averages this value over publications. Hit@1 is the fraction of publications
whose first result is recorded as relevant. Micro Recall@100 divides the total
recorded links recovered in the first 100 results by all recorded links across
publications. Equal logits retain catalog order. These definitions and the
acceptance tolerance $5\times10^{-5}$ are implemented in
\evidence{src/reaction_lens/evaluation.py}{the evaluation script}.

This exercise verifies packaging and cached-feature inference. It does not
independently reproduce encoder training, head training or every historical
experiment. The complete-text latency experiment exercises the local encoder but
does not rebuild the entire development feature cache. The reserved test remains
unopened. The measured scorer, encoder, pipeline and benchmark sources are preserved by
checksums; packaging, validation and documentation can be revised separately.

\section{Limitations and practical implications}
This is a single-system, single-checkpoint, single-host experience report. It does
not demonstrate a general improvement over vector databases, other retrieval
architectures or other testing methods. Candidate independence follows directly
from pointwise scoring; the evidence is about an implementation's conformance and
costs, not discovery of a new mathematical property. We have not measured failures
in deployed systems, mutation-detection rates, user productivity or maintenance
effort. The contract checks contain related scalar comparisons, not independent
statistical samples. The convenience text sample and repeated development use
limit semantic and performance generalization.

The reusable engineering lessons are to state whether a score is pairwise or
candidate-set-dependent, distinguish score stability from rank stability, test
blocking separately from retrieval quality, and report encoder and head costs
separately. For reproduction, pin the model and feature artifacts and distinguish
reconstruction of tables from actual inference. These practices can make this
system easier to inspect, but our study does not quantify developer benefits.

\section{Conclusion and availability}
Reaction Lens Systems provides a concrete case in which exact catalog scoring can be
blocked without observed changes to retained candidates' scores, while encoder
execution dominates measured local request latency. The artifact records all
probes and ties displayed results to raw measurements through a generated
\texttt{NUMBERS.csv}. The accompanying instructions provide table reconstruction,
fresh inference, and experiment rerun commands as distinct operations.
The code is licensed under MIT and the manuscript under CC BY 4.0.

\paragraph{Tool assistance.} An AI coding assistant assisted with the measurement
scripts, analysis and manuscript preparation. Reported measurements were produced
by executing the scripts. The author is responsible for the content.

{\small
\setlength{\bibsep}{3pt}
\bibliographystyle{plainnat}
\bibliography{references}}
\end{document}
